% WORK-SUMMARY.tex
% Detailed summary of all work contained in the `nnip/` repository.
% Compile with: pdflatex WORK-SUMMARY.tex  (run twice for the table of contents)
\documentclass[11pt,a4paper]{article}

\usepackage[T1]{fontenc}
\usepackage[utf8]{inputenc}
\usepackage[margin=2.3cm]{geometry}
\usepackage{amsmath,amssymb}
\usepackage{booktabs}
\usepackage{longtable}
\usepackage{array}
\usepackage{enumitem}
\usepackage{xcolor}
\usepackage[hidelinks]{hyperref}

\newcommand{\code}[1]{\texttt{\small #1}}
\newcommand{\pct}{\,\%}

% long unbreakable paths inside \code: prefer loose lines over margin overflow
\setlength{\emergencystretch}{3em}
\sloppy

\title{\textbf{Maxwell-Fluid Shrinking-Sheet PINN:\\
Reproduction, Evaluator Harness, GPU Candidate Search\\
and Stability-Table Audit}\\[0.6em]
\large Complete record of the work in the \code{nnip/} repository}
\author{Raja Rawat \quad (work carried out with Claude Code / Open Discovery harness)}
\date{Summary compiled 12 November 2026 \\ \small Work performed 5--6 November 2026 (three sessions)}

\begin{document}
\maketitle

\begin{abstract}
\noindent
This document is a complete, self-contained summary of the work stored in the
\code{nnip/} repository. The subject is a Physics-Informed Neural Network (PINN)
reproduction of Pandey \textit{et al.}, \emph{``Insight of thermal transport in Maxwell
fluid flow on a porous contracting flat surface inside a porous material with
non-uniform heat source/sink: Multiple solutions and stability treatment''},
\emph{Forces in Mechanics} 100352 (DOI 10.1016/j.finmec.2026.100352).

Three things were produced. \textbf{(i) A reproduction harness:} the paper's equations,
boundary conditions and Tables~1--3 were transcribed into a durable provenance record,
and a candidate-proof \emph{external evaluator} was written that gates a PINN script
against the paper's published numbers \emph{and} two independent numerical oracles. The
baseline PINN was hash-locked twice (CPU and GPU) through the Open Discovery
\code{program-evolution} contract. \textbf{(ii) A GPU candidate search:} five candidates
were proposed and evaluated under a pre-registered acceptance rule; the winner
(candidate-005, a combination of four orthogonal mechanisms) is
\textbf{2.72$\times$ faster end-to-end on the same GPU} with accuracy unchanged at the
search point, and was subsequently shown \textbf{accuracy-equivalent} to the baseline on
unseen parameter rows and seeds. \textbf{(iii) A stability-table audit:} Tables~1 and~2
and the critical suction parameter $S_c$ reproduce; \textbf{Table~3's eigenvalue
magnitudes do not}. Three independent solvers of the paper's own Eq.~(31) agree with each
other to $<10^{-7}$ relative and disagree with the published lower-branch column by factors
of 2.8--13$\times$, while the \emph{signs} (upper stable, lower unstable) do reproduce. A
plausible mechanical explanation for the published upper column is identified: the bundled
Chebyshev oracle applies its boundary rows at swapped ends, and that placement reproduces
all five published upper values to within $0.7\pct$. Section~\ref{sec:beat} weighs the
whole record against the article itself: superseded on the stability table, improved on the
far-field reference, matched on Tables~1 and~2, and still well behind on computational cost
and parameter coverage.
\end{abstract}

\tableofcontents
\clearpage

%=====================================================================
\section{Scope, inputs and provenance}

\subsection{What the human supplied}
A single chat session provided four artifacts, all originating outside the repository:

\begin{enumerate}[leftmargin=1.4em,itemsep=2pt]
  \item The paper PDF (12 pages, uncorrected proof), \code{pagination\_FINMEC\_100352-1.pdf}.
  \item \code{pinn\_maxwell\_shrinking.py} --- a three-stage PINN solving the paper's
        similarity ODEs (branch solve, fold tracing, stability eigenvalue).
  \item \code{reference\_bvp.py} --- an independent \code{scipy.solve\_bvp} collocation
        solver for the same boundary-value problem, used as ground truth.
  \item \code{ref\_eig\_cheb.py} --- an independent Chebyshev-spectral solve of the
        linearised stability eigenvalue problem, a second ground truth.
\end{enumerate}

\noindent The instruction was to use the existing Open Discovery harness
(\code{paper-implementer} for the reproduction claim, \code{program-evolution} for an
optimisation baseline) to give the work a durable, evidence-preserving home.

\subsection{Non-destructive provenance}
Per the harness contract, external work is never modified. The scripts and PDF were copied
read-only into the project and their SHA-256 digests recorded (Table~\ref{tab:prov}).
No upstream Open Discovery file was edited: a check of the working tree against the
upstream history (\code{open-discovery/.git-upstream}) shows \emph{zero} modified tracked
files. All new material lives in two directories that the harness \code{.gitignore}
deliberately excludes and that were force-added to the outer repository:
\code{initiatives/maxwell-shrinking-pinn/} and \code{skill-lab/evolve-maxwell-pinn-gpu/}.

\begin{table}[ht]
\centering
\caption{Source-material SHA-256 digests recorded at copy time (truncated).}
\label{tab:prov}
\small
\begin{tabular}{@{}ll@{}}
\toprule
\textbf{File} & \textbf{SHA-256 (prefix)} \\
\midrule
\code{pinn\_maxwell\_shrinking.py} & \code{8f145d68377b7be4bbbc73f\ldots b180} \\
\code{reference\_bvp.py}           & \code{a1e042fd16bdccd19cdac98\ldots f203} \\
\code{ref\_eig\_cheb.py}           & \code{081318bf7cbb5ab4a78d7e9\ldots 60b1} \\
\code{pagination\_FINMEC\_100352-1.pdf} & \code{005acd8e27c6533be8e9e4e\ldots f5e3} \\
\midrule
Locked baseline bundle (\code{lock.json} root) & \code{40d7e8611b8d417153eea25\ldots c281} \\
Evaluator (\code{evaluator.py})               & \code{e99f0be74b50d55d74e5e8a\ldots 40eb} \\
Candidate-005 winner, rebuilt from patch      & \code{cae54498\ldots 0661} \\
\bottomrule
\end{tabular}
\end{table}

\subsection{Repository layout}
\begin{itemize}[leftmargin=1.4em,itemsep=1pt]
  \item \code{nnip/amit-pinn-article/} --- the original source material (PDF, three
        scripts) plus \code{requirements-gpu-cu126.txt}, the exact pinned GPU environment
        (torch 2.14.1+cu126, scipy 1.18.1, numpy 2.5.4). The 6.7\,GB venv itself is
        excluded from version control and rebuilt from this file.
  \item \code{\ldots/open-discovery/} --- the research harness (upstream, unmodified):
        \code{AGENTS.md}, 20 workflow skills under \code{.agents/skills/}, seven agent
        role definitions, \code{docs/} (evidence standard, candidate lifecycle,
        state model), and the \code{program-evolution/} locking harness.
  \item \code{\ldots/initiatives/maxwell-shrinking-pinn/} --- \textbf{all new scientific
        work} (2.2\,MB): \code{BRIEF.md} and the single project
        \code{projects/pinn-reproduction/}.
  \item \code{\ldots/skill-lab/evolve-maxwell-pinn-gpu/} --- \textbf{new} prototype skill:
        how to run a GPU candidate search for this specific initiative, plus a menu of
        concrete candidate mechanisms.
\end{itemize}

%=====================================================================
\section{The physical problem and the claims under test}

\subsection{Governing similarity ODEs (paper Eqs.~13--15)}
Two-dimensional steady boundary-layer flow and heat transfer of a Maxwell (viscoelastic)
fluid over a continuously \emph{contracting} sheet embedded in a Darcy porous medium, with
wall mass transfer and a heat source/sink depending on \emph{both} temperature and spatial
position:
\begin{align}
  (1-\beta f^2)\,f''' + 2\beta f f' f'' + f f'' - f'^2 - K f' &= 0, \label{eq:mom}\\
  \theta'' + \mathrm{Pr}\, f\,\theta' + A^{*}\theta + B^{*}e^{-\eta} &= 0, \label{eq:en}
\end{align}
subject to
\begin{equation}
  f(0)=S,\quad f'(0)=-1,\quad f'(\infty)=0,\qquad \theta(0)=1,\quad \theta(\infty)=0 .
\end{equation}
Here $\beta$ is the Deborah (Maxwell) number, $K$ the permeability parameter, $S$ the wall
mass-transfer parameter ($S>0$ suction), $\mathrm{Pr}$ the Prandtl number, and
$A^{*},B^{*}$ the temperature- and space-dependent source/sink strengths. The paper
truncates the far field at $\eta_\infty=10$. The engineering outputs are
\begin{equation}
  C_f\,\mathrm{Re}_x^{1/2} = (1+\beta) f''(0), \qquad
  \mathrm{Nu}\,\mathrm{Re}_x^{-1/2} = -\theta'(0).
\end{equation}

\subsection{The paper's own method, and why cross-validation matters}
The paper solves the system by shooting (RK4, $\Delta\eta=0.01$, Secant iteration to
$10^{-5}$). Neither companion script uses that route: \code{reference\_bvp.py} uses
collocation and the PINN uses a neural trial function with hard-constrained boundary
conditions. Three independent numerical routes to the same claim is exactly the right
design for an audit, and this independence is what later made the Table~3 discrepancy
diagnosable rather than ambiguous.

\subsection{Claims transcribed as automated gates}
\code{paper/paper-summary.md} transcribes the equations, boundary conditions, solution
method, all of Tables~1--3, and the qualitative conclusions of Section~6. Each number is
explicitly marked either \textbf{[gate]} (used as an automated evaluator gate,
cross-checked against the co-author's own inline comments in the original script) or
\textbf{[reference only --- verify against PDF before quoting]} (read visually off rendered
table images). The gated values are listed in Table~\ref{tab:gates}.

\begin{table}[ht]
\centering
\caption{The published values used as automated evaluator gates. Baseline parameter row:
$\beta=0.1$, $K=0.1$, $S=1.95$, $\mathrm{Pr}=2$, $A^{*}=B^{*}=0.1$; stability row $S=2.0$,
$B^{*}=0$.}
\label{tab:gates}
\small
\begin{tabular}{@{}llrr@{}}
\toprule
\textbf{Claim} & \textbf{Source} & \textbf{Upper branch} & \textbf{Lower branch} \\
\midrule
$C_f \mathrm{Re}^{1/2}$          & Table 1, p.~7  & 1.93523180 & 0.75244081 \\
$\mathrm{Nu}\,\mathrm{Re}^{-1/2}$ & Table 2, p.~7  & 3.35807824 & 3.14427924 \\
Critical suction $S_c$            & Sec.~4.1, Fig.~2 & \multicolumn{2}{c}{1.8391} \\
Leading eigenvalue $\alpha_1$     & Table 3, p.~10 & $+0.8810518$ & $-0.0376726$ \\
\bottomrule
\end{tabular}
\end{table}

\noindent The physics being protected: dual (upper/lower) similarity solutions exist only
for $S\geq S_c$; the upper branch is stable ($\alpha_1>0$) and the lower branch unstable
($\alpha_1<0$) at every tabulated $S$. A ``faster'' candidate that converges to one branch,
swaps branch labels, or flips an eigenvalue sign has improved nothing, however low its
wall-clock.

%=====================================================================
\section{The baseline program}

\code{baseline/pinn\_maxwell\_shrinking.py} (290 lines, float64 throughout, device
auto-selected as CUDA when available) implements three separate PINNs:

\begin{enumerate}[leftmargin=1.4em,itemsep=2pt]
  \item \code{solve\_branch(p, branch, \ldots)} --- the upper or lower branch at one
        parameter point. A $4\times40$ tanh MLP (Xavier-normal init, zero bias) on the
        mapped variable $\xi = 1-e^{-\eta/L}$ with 300 collocation points up to
        $\xi_{\max}=0.9995$. Far-field boundary conditions are imposed \emph{exactly} by
        the trial function, not by penalty. Training is 4000 Adam
        ($\mathrm{lr}=2\times10^{-3}$) $+$ 1500 LBFGS for momentum, then
        $2000+1000$ for energy, preceded by a $2500+200$ ``pilot'' solve whose only purpose
        is to read off $f(\infty)$ and choose the map length $L_{\mathrm{MAP}}$
        (\code{auto\_scale=True}). Branch selection uses an \code{f2\_anchor} warm start.
  \item \code{trace\_fold(p, f2\_values)} --- treats $S$ as the unknown and prescribes
        $f''(0)$ over eight continuation values, so the solution path passes smoothly
        through the turning point and yields $S_c$ directly.
  \item \code{smallest\_eigen(base\_net, p, alpha\_init)} --- the leading eigenvalue
        $\alpha_1$ of the linearised disturbance problem about a frozen base flow
        ($3000+1000$ iterations, $\alpha$ a trainable parameter).
\end{enumerate}

\noindent The script's own docstring records that $L_{\mathrm{MAP}}$ auto-scaling is what
brings the lower branch from a $1$--$2\pct$ error floor down to ${\sim}0.1\pct$; this is
flagged in the project notes as something a candidate must not silently remove.

%=====================================================================
\section{Infrastructure built}

\subsection{The external evaluator (\code{evaluator.py}, 364 lines)}
A standalone, \emph{import-based} (not print-parsing) evaluator implementing the
\code{program-evolution} contract: it exits 0 and prints exactly one JSON object. It loads
a candidate module via \code{importlib} and first checks the required public contract
\code{\{MLP, train, fields, wall\_values, solve\_branch, trace\_fold, smallest\_eigen\}}.
A candidate may change anything \emph{inside} those functions; it may not change their
signatures, and it may never edit the evaluator.

\begin{description}[leftmargin=1.6em,itemsep=2pt]
  \item[\code{--mode smoke} (default, cheap).] Monkeypatches the candidate's \code{train}
        to a hard 120-Adam/40-LBFGS budget and checks the dual-solution pipeline runs and
        returns finite, physically ordered values ($C_f^{\mathrm{up}} > C_f^{\mathrm{low}} > 0$,
        both $\mathrm{Nu}>0$). This is a wiring/pipeline gate, explicitly \emph{not} an
        accuracy claim --- the evaluator docstring says so, and the project notes repeat it.
  \item[\code{--mode full} (expensive).] Runs the candidate's own unmodified
        \code{solve\_branch} at the paper's baseline row and gates on relative error to the
        published $C_f$ and $\mathrm{Nu}$ (default tolerance $1\pct$).
  \item[\code{--fold}] gates $S_c$ against 1.8391 (tolerance $3\pct$ in the acceptance rule).
  \item[\code{--stability}] gates the Table~3 eigenvalues, including a \emph{hard sign
        check} (upper $>0$, lower $<0$) --- the real guard against branch confusion.
  \item[\code{--oracle}] cross-checks the \emph{paper's} numbers against the independent
        \code{scipy.solve\_bvp} oracle. Informative only; it never gates the candidate.
\end{description}

\noindent The objective reported is \code{elapsed\_seconds} (minimised), summed over the
stages actually requested.

\subsection{Baseline locks}
\code{program-evolution/scripts/lock\_baseline.py} snapshots baseline $+$ evaluator with
hashes, environment, and a smoke run.

\begin{table}[ht]
\centering
\caption{The two baseline locks. Both report \code{status: locked}, \code{error: null},
\code{valid: true}, \code{SUPPORTED}.}
\label{tab:locks}
\small
\begin{tabular}{@{}lll@{}}
\toprule
 & \code{baseline-lock/} (session 1) & \code{baseline-lock-gpu/} (session 2) \\
\midrule
Hardware     & CPU only (Windows 10) & NVIDIA GTX 1080 Ti (sm\_61) \\
torch        & 2.13.0+cpu & 2.14.1+cu126 (driver 580.178.04) \\
Python / OS  & 3.13.7 / Windows-10 & 3.12.15 / Linux 6.8 glibc 2.39 \\
Smoke objective & \textbf{126.81\,s} & \textbf{14.77\,s} \\
Smoke $C_f$ upper / lower & 1.9125139387 / 0.7721875369 & 1.9125139387 / 0.7721875369 \\
\bottomrule
\end{tabular}
\end{table}

\noindent Note that the two locks reproduce the smoke values to 10--11 significant figures
across operating system, Python version, torch version and device --- a strong
cross-platform determinism check that was obtained for free from the locking discipline.

\subsection{The GPU-search skill and mechanism menu}
\code{skill-lab/evolve-maxwell-pinn-gpu/SKILL.md} is a new repository-local skill. Its
substantive content is methodological discipline:
\begin{enumerate}[leftmargin=1.4em,itemsep=1pt]
  \item Verify CUDA is genuinely present before any speed claim.
  \item \textbf{Never compare wall-clock across hardware.} Moving to a GPU is a hardware
        effect, not an evolved program; a fresh same-hardware reference must be established
        first, and the hardware and algorithmic improvements must be reported separately
        and \emph{never summed}.
  \item One candidate folder per hypothesis, gate on smoke first, spend full-mode budget
        only on survivors, never edit the evaluator.
  \item A claimed winner must be \emph{reconstructed} from the locked baseline plus its
        recorded patch and re-confirmed on fresh parameter points over several seeds.
  \item Novel-science extensions are a \emph{new project}, not a candidate.
\end{enumerate}
\code{references/candidate-mechanisms.md} is a menu of 14 concrete, line-referenced
hypotheses in four groups: \textbf{A} same-hardware throughput (A1 batch both branches via
\code{vmap}; A2 cheaper $L_{\mathrm{MAP}}$ pilot; A3 early-exit LBFGS; A4 Adam/LBFGS
rebalance; A5 vectorise the fold continuation loop; A6 trade collocation count for
iteration count), \textbf{B} accuracy and robustness (closed-form $L_{\mathrm{MAP}}$, more
reliable branch selection, stress-testing the $(1-\beta f^2)$ singularity, targeting the
lower branch), \textbf{C} GPU-native restructuring under the same public contract, and
\textbf{D} a novel-work track (parametric/meta-PINN over $(S,\beta,K)$, continuous
fold-tracking by differentiating with respect to $S$, parametric stability locating
$\alpha_1=0$ as an independent check that the turning point and the stability-exchange
point coincide, and ensemble uncertainty).

\subsection{Independent confirmation and diagnostic tooling (new code)}
\begin{itemize}[leftmargin=1.4em,itemsep=1pt]
  \item \code{confirm\_fresh.py} (55 lines) --- runs a program's own \code{solve\_branch}
        at three Table~1 rows \emph{not} used during the search
        ($S{=}2.10$; $K{=}0.13$; $\beta{=}0.15$) over several seeds, and compares $C_f$ and
        $\mathrm{Nu}$ both with the paper and with a long-domain
        ($\eta_\infty{=}40$, converged) \code{scipy.solve\_bvp} oracle.
  \item \code{diagnostics/eig\_shooting.py} --- independent IVP shooting solve of the
        paper's Eq.~(31), with a domain-length study.
  \item \code{diagnostics/eig\_cheb\_fixed.py} --- the bundled Chebyshev oracle with its
        boundary rows moved to the correct ends.
  \item \code{oracle-v2/table3\_oracle.py} (139 lines) $+$ \code{judge.py} --- a
        full Table~3 oracle (three methods $\times$ five domain lengths $\times$ five $S$
        values $\times$ two branches) with the pre-registered acceptance criteria applied
        mechanically to produce \code{judge.md}.
  \item \code{summarize\_run.py}, \code{run\_candidates.sh}, \code{queue\_round2.sh},
        per-run \code{run.sh}/\code{timeline.txt}/\code{env.txt}/\code{nvidia-smi.txt}/
        \code{sha256.txt} --- the evidence plumbing.
\end{itemize}

%=====================================================================
\section{Session 1 (5 Nov 2026, CPU): harness bring-up}

\begin{itemize}[leftmargin=1.4em,itemsep=1pt]
  \item Paper read in full (12 pages); equations, boundary conditions, solution method and
        Tables~1--3 transcribed with gate/reference-only marking.
  \item Source material copied read-only with hashes recorded; originals untouched.
  \item \code{evaluator.py} written and validated in smoke mode against the baseline:
        \textbf{valid, SUPPORTED}, $C_f^{\mathrm{up}}{\approx}1.913$,
        $C_f^{\mathrm{low}}{\approx}0.772$, $\mathrm{Nu}^{\mathrm{up}}{\approx}1.854$,
        $\mathrm{Nu}^{\mathrm{low}}{\approx}1.203$ in 126.8\,s --- correctly ordered and
        finite, and explicitly \emph{not} claimed to match the paper.
  \item Baseline locked (\code{baseline-lock/lock.json}, \code{status: locked}).
  \item A full-mode CPU confirmation run was launched in the background. \textbf{It never
        completed} --- \code{runs/full\_table12\_confirmation.json} is empty, and this is
        recorded as such rather than quietly dropped. The notes stated in advance that no
        pass/fail may be reported before the process exits.
\end{itemize}

%=====================================================================
\section{Session 2 (5 Nov 2026, GTX 1080 Ti): GPU reference and candidate search}

\subsection{Pre-registration, written before any candidate result}
\code{runs/GPU-SEARCH-PREREGISTRATION.md} fixes the acceptance rule \emph{before} results
exist:
\begin{enumerate}[leftmargin=1.4em,itemsep=1pt]
  \item smoke gate valid;
  \item full mode: no Table~1/2 failure (tol $1\pct$), no $S_c$ failure (tol $3\pct$);
  \item stability: both sign checks pass; $\alpha_1$ magnitudes \emph{reported} against the
        independent values, not used to accept or reject;
  \item the speed claim is the same-GPU total objective versus 1451.8\,s;
  \item a winner must be rebuilt from baseline $+$ recorded patch and re-confirmed on fresh
        parameter points over several seeds.
\end{enumerate}

\noindent Two reasons are documented for excluding the $\alpha_1$ magnitude gate, both
established before candidate results:

\begin{itemize}[leftmargin=1.4em,itemsep=2pt]
  \item \textbf{No correct solver can satisfy it.} Two independent solves of the paper's own
        Eq.~(31), verified term-by-term against page~10, agree with each other and not with
        Table~3: shooting gives $\alpha_1^{\mathrm{up}}=0.7358$,
        $\alpha_1^{\mathrm{low}}=-0.4938$ at $S{=}2.0$; the corrected Chebyshev solve gives
        identical roots. $|\alpha_1^{\mathrm{low}} + 0.0377| < 1\pct$ is unreachable. The
        evaluator was \emph{not} edited --- the program-evolution contract forbids weakening
        a gate mid-search; the gate is simply reported.
  \item \textbf{The upper-branch magnitude is not identifiable by this PINN.} In the far
        field $g'\sim e^{-\mu\eta}$ with
        $(1-\beta f_\infty^2)\mu^2 - f_\infty\mu + (\alpha-K) = 0$. For $\alpha > K$ both
        roots have $\mathrm{Re}\,\mu>0$, so $g'(\infty)=0$ holds for a \emph{continuum} of
        $\alpha$ on the semi-infinite domain. The PINN trial function imposes
        $g'(\infty)=0$ exactly, so the upper root ($\approx 0.74 > K = 0.1$) is not pinned
        --- CPU tests with initialisations spanning $0.0$--$1.0$ all land in $0.66$--$0.69$,
        while the GPU baseline lands at $0.958$. The lower branch
        ($\alpha\approx-0.49 < K$) has one growing mode, hence a genuine discrete
        eigenvalue, and the PINN matches the independent value to ${\sim}1\pct$. Conclusion
        recorded in advance: \emph{only the upper sign is a meaningful PINN output}, and no
        candidate targets the upper magnitude.
\end{itemize}

\subsection{The same-GPU reference run}
\code{runs/full\_confirmation\_gpu.json}: unmodified baseline, unmodified evaluator,
\code{--mode full --oracle --fold --stability}. Total objective \textbf{1451.79\,s}
(Table~1/2 404.54\,s, fold 456.02\,s, stability 591.23\,s).

\begin{table}[ht]
\centering
\caption{GPU reference run against the paper. Table~1/2 and $S_c$ pass at $1\pct$ and
$3\pct$; Table~3 magnitudes fail; both signs pass.}
\label{tab:gpuref}
\small
\begin{tabular}{@{}llrr@{}}
\toprule
\textbf{Quantity} & \textbf{Paper} & \textbf{PINN (GPU)} & \textbf{Rel.\ error} \\
\midrule
$C_f$ upper        & 1.93523180 & 1.93516109 & $0.0037\pct$ \\
$C_f$ lower        & 0.75244081 & 0.75562700 & $0.423\pct$ \\
$\mathrm{Nu}$ upper & 3.35807824 & 3.35003942 & $0.239\pct$ \\
$\mathrm{Nu}$ lower & 3.14427924 & 3.13232975 & $0.380\pct$ \\
$S_c$              & 1.8391     & 1.84910    & $0.543\pct$ \\
$\alpha_1$ upper   & $+0.8810518$ & $+0.958033$ & $8.74\pct$ \;(\textbf{FAIL}) \\
$\alpha_1$ lower   & $-0.0376726$ & $-0.488198$ & $1195.9\pct$ \;(\textbf{FAIL}) \\
\bottomrule
\end{tabular}
\end{table}

\noindent The \code{--oracle} cross-check confirms the transcription and the paper's own
Table~1/2 numbers independently: \code{scipy.solve\_bvp} returns
$C_f^{\mathrm{up}}=1.93525922$ ($1.4\times10^{-3}\pct$ from the paper),
$C_f^{\mathrm{low}}=0.75243975$ ($1.4\times10^{-4}\pct$),
$\mathrm{Nu}^{\mathrm{up}}=3.35133372$ ($0.20\pct$),
$\mathrm{Nu}^{\mathrm{low}}=3.13525713$ ($0.287\pct$). In other words the oracle is itself
$0.2$--$0.3\pct$ from the paper's $\mathrm{Nu}$ values, so the PINN's $0.24$--$0.38\pct$
is at the resolution of the reference, not a PINN deficiency.

\subsection{Candidate lineage}
Five candidates, each a full copy of the baseline with one recorded change
(\code{patch.diff}), one hypothesis (\code{NOTE.md}), and its own evaluator outputs. All
five passed smoke (8.0--11.9\,s) and all five passed the pre-registered full-mode gates.

\begin{table}[ht]
\centering
\caption{Candidate lineage, single runs, same GTX 1080 Ti, unmodified evaluator.
The ``vs baseline'' column is against the 1451.8\,s reference.}
\label{tab:lineage}
\small
\begin{tabular}{@{}llp{5.4cm}rrl@{}}
\toprule
\textbf{ID} & \textbf{Parent} & \textbf{Change} & \textbf{Total (s)} & \textbf{Speed-up} & \textbf{Gates} \\
\midrule
001 & baseline & Taylor-mode derivative jets instead of nested autograd & 823.5 & 1.76$\times$ & pass \\
002 & baseline & adaptive LBFGS stop (50-iter chunks, stop at $<1\pct$ gain) & 967.7 & 1.50$\times$ & pass \\
003 & baseline & collocation $\xi_{\max}$ $0.9995 \to 0.99995$ & 1374.9 & none (same work) & pass\footnotemark[1] \\
004 & baseline & auto-scale pilot $2500/200 \to 800/100$ & 1230.7 & ${\sim}1.18\times$\footnotemark[2] & pass \\
\textbf{005} & 001 (+\,2/3/4) & combination of all four mechanisms & \textbf{503.5} & \textbf{2.88$\times$} & pass \\
\bottomrule
\end{tabular}
\end{table}
\footnotetext[1]{Candidate-003's CPU accuracy gain was \emph{not} reproduced in full mode --- recorded as a negative result rather than dropped.}
\footnotetext[2]{Within run-to-run noise; single-run stage timings vary ${\sim}15\pct$ (the fold stage, untouched by 003/004, ran 368--387\,s against the reference's 456\,s).}

\paragraph{Candidate-001 (the main mechanism).} Every optimiser step of the baseline builds
$f',f'',f'''$ (and $\theta',\theta''$, $g',g'',g'''$) with three nested
\code{torch.autograd.grad(create\_graph=True)} passes and then backpropagates through that
nested graph. For a $4\times40$ tanh MLP on 300 points this is almost pure kernel-launch
overhead on a GPU. Candidate-001 propagates the value and its first $k$ $\eta$-derivatives
forward in closed form (Fa\`a di Bruno through \code{tanh}, Leibniz for the trial-function
products), giving an identical residual with a much smaller graph. Architecture,
initialisation, seeds, collocation, iteration counts and optimiser settings are unchanged.
Verification before any evaluator run: on random weights the jet residuals \emph{and their
parameter gradients} match baseline autograd to $\leq 2\times10^{-14}$ in float64, and the
eigen trial-function jet to $\leq 1.1\times10^{-15}$.

\paragraph{Candidate-002.} The LBFGS tolerances ($10^{-12}$ gradient, $10^{-15}$ change)
never trigger, so every stage burns its full fixed budget. A CPU trace showed momentum
converges in ${\sim}100$ LBFGS evaluations and energy in ${\sim}200$; the remaining
${\sim}1700$ and ${\sim}1050$ evaluations move $C_f$/$\mathrm{Nu}$ by $<10^{-5}$ relative.
The candidate runs the same LBFGS in 50-iteration chunks (history preserved) and stops when
a chunk improves the loss by $<1\pct$, never exceeding the original budget.

\paragraph{Candidate-003.} $\xi_{\max}=0.9995$ enforces the residual only out to
$\eta \approx 7.6L$, while the lower branch's slow decay ($\lambda\approx0.77$) still
matters to $\eta\approx20$. Raising $\xi_{\max}$ to $0.99995$ extends the reach to
${\sim}13L$. A CPU pre-check showed lower-$C_f$ error against the $\eta_\infty\to\infty$
oracle falling from $0.19\pct$ to $0.018\pct$ --- but this gain did \emph{not} reproduce in
full mode, which is recorded explicitly.

\paragraph{Candidate-004.} The pilot only needs $f(\infty)$ to choose
$L=\lceil\lambda\rceil/\lambda$. A budget sweep ($2500/200$, $1500/100$, $800/100$,
$400/60$) moved the chosen lower-branch $L$ only within $1.3095$--$1.3111$, so the pilot was
cut to $800+100$.

\paragraph{Candidate-005.} The four mechanisms act on different things (derivative
evaluation, LBFGS stopping, collocation reach, pilot budget), so the hypothesis was that the
gains multiply. Candidate-005 is candidate-001 plus the recorded baseline-relative patches
of 002, 003 and 004 applied in order; all hunks applied cleanly. The note records the
discipline that it is only to be trusted because each component passed independently in
round~1.

\subsection{Winner confirmation}
The winner was rebuilt from \code{baseline-lock-gpu/baseline/} plus
\code{candidate-005/patch.diff} and verified byte-identical to the evaluated candidate
(SHA-256 \code{cae54498\ldots 0661}). Paired full-mode runs were then taken back-to-back on
the same otherwise-idle GPU (17:49--18:44 UTC, timeline recorded).

\begin{table}[ht]
\centering
\caption{Paired same-GPU full-mode runs (\code{--oracle --fold --stability}), unmodified
evaluator. \textbf{Algorithmic improvement: 2.72$\times$.}}
\label{tab:paired}
\small
\begin{tabular}{@{}lrrr@{}}
\toprule
\textbf{Stage} & \textbf{Baseline (s)} & \textbf{Candidate-005 (s)} & \textbf{Speed-up} \\
\midrule
Table 1/2 (both branches, incl.\ pilots) & 392.8 & 158.1 & 2.48$\times$ \\
Fold / $S_c$                             & 396.4 & 161.9 & 2.45$\times$ \\
Stability (2 branch solves + 2 eigen)    & 600.7 & 191.7 & 3.13$\times$ \\
\midrule
\textbf{Total}                           & \textbf{1389.9} & \textbf{511.6} & \textbf{2.72$\times$} \\
\bottomrule
\end{tabular}
\end{table}

\noindent The hardware effect is reported separately and \emph{not} summed: the Windows-CPU
smoke lock took 126.8\,s against the GPU smoke lock's 14.8\,s --- different machines, smoke
mode only.

\begin{table}[ht]
\centering
\caption{Accuracy at the search point is unchanged within noise.}
\label{tab:acc005}
\small
\begin{tabular}{@{}lllp{5.2cm}@{}}
\toprule
\textbf{Metric} & \textbf{Baseline} & \textbf{Candidate-005} & \textbf{Reference} \\
\midrule
$C_f$ upper & 1.935161 ($0.004\pct$) & 1.934442 ($0.041\pct$) & paper 1.93523180 \\
$C_f$ lower & 0.755627 ($0.42\pct$)  & 0.755643 ($0.43\pct$)  & paper 0.75244081; converged oracle 0.753702 \\
$\mathrm{Nu}$ upper / lower error & $0.24\pct$ / $0.38\pct$ & $0.27\pct$ / $0.39\pct$ & the oracle is itself $0.20\pct$ / $0.29\pct$ from the paper \\
$S_c$ & 1.84910 ($0.54\pct$) & 1.84947 ($0.56\pct$) & paper 1.8391 \\
$\alpha_1$ lower & $-0.48820$ & $-0.48808$ & independent $-0.4938$ (paper $-0.0377$) \\
$\alpha_1$ upper (sign only) & $0.958$ & $1.051$ & independent leading $0.7358$ (paper $0.881$) \\
\bottomrule
\end{tabular}
\end{table}

\subsection{Fresh-parameter check and the honest open item}
\code{confirm\_fresh.py} was run on candidate-005 at three unseen Table~1 rows over seeds
1--3 (the search used seed~0 and only the $S{=}1.95$ and $S{=}2.0$ rows). Errors against
the converged $\eta_\infty{=}40$ oracle: upper $C_f$ $0.03$--$0.21\pct$, all
$\mathrm{Nu}$ $\leq0.34\pct$, lower $C_f$ $0.04$--$1.00\pct$. Against the paper's own
Table~1 $C_f$, 16 of 18 points fall within $1\pct$; the two exceptions are the lower branch
at $S{=}2.1$ (seeds 2 and 3: $1.51\pct$, $1.79\pct$), where \emph{the paper's own value} is
$0.78\pct$ below the converged oracle --- an $\eta_\infty{=}10$ truncation effect that grows
as the lower branch decays more slowly.

The session-2 verdict deliberately left one item open rather than overclaiming: whether the
lower-branch seed spread is inherent to the baseline architecture or was \emph{introduced}
by candidate-005 was untested, because the standard plan had not run the baseline through
\code{confirm\_fresh.py}. That comparison was named as the required next step before calling
005 accuracy-equivalent off the search point.

%=====================================================================
\section{Session 3 (6 Nov 2026, RTX 2060 SUPER): closing both open items}

The venv was rebuilt from \code{requirements-gpu-cu126.txt} (torch 2.14.1+cu126, scipy
1.18.1 --- identical to session~2). Because the GPU differs, \emph{no} timing from this
session is compared with session-2 timings. Both items were pre-registered in
\code{runs/OPEN-ITEMS-PREREGISTRATION.md} with program hashes verified before running.

\subsection{Item A --- is the lower-branch seed spread candidate-005's fault?}
\textbf{Pre-registered rule.} Run unmodified \code{confirm\_fresh.py} on both programs,
seeds 1--3, serially on the same idle GPU (baseline first); 9 (row, seed) points per
program. Primary metric: lower-branch $C_f$ error against the $\eta_\infty{=}40$ oracle.
With $m$ and $x$ the mean and max over the nine points, declare \emph{equivalent} if
$|m_5-m_B|\leq0.20$ percentage points \emph{and} $x_5\leq x_B+0.30$\,pp; \emph{005 worse} or
\emph{005 better} otherwise. The pre-registration states plainly that $n=9$ per arm detects
only gross differences, that ``equivalent'' means ``no gross difference detected'' and not
proof of equality, and that the two programs consume the RNG differently so the points are
not truly paired.

\textbf{Result: equivalent} on the primary metric and all three secondaries
(Table~\ref{tab:itemA}). The lower-branch $C_f$ seed spread ($0.03$--$1.0\pct$) is present
in the \emph{untouched baseline}; candidate-005 did not introduce it.

\begin{table}[ht]
\centering
\caption{Item A, errors against the converged oracle over 9 (row, seed) points each,
RTX 2060 SUPER, job pinned to 4 CPU cores, 04:54--08:12 UTC.}
\label{tab:itemA}
\small
\begin{tabular}{@{}lrrl@{}}
\toprule
\textbf{Metric} & \textbf{Baseline mean / max ($\pct$)} & \textbf{005 mean / max ($\pct$)} & \textbf{Verdict} \\
\midrule
Lower $C_f$ (\textsc{primary}) & 0.459 / 0.962 & 0.392 / 0.988 & equivalent \\
Upper $C_f$                    & 0.172 / 0.304 & 0.101 / 0.223 & equivalent \\
Upper $\mathrm{Nu}$            & 0.058 / 0.156 & 0.090 / 0.183 & equivalent \\
Lower $\mathrm{Nu}$            & 0.147 / 0.333 & 0.180 / 0.347 & equivalent \\
\bottomrule
\end{tabular}
\end{table}

\paragraph{Side findings (informative).} Candidate-005 on this GPU reproduces its
session-2 (1080 Ti) values to $\leq0.05$\,pp at 8 of 9 points; the ninth
($S{=}1.95,\beta{=}0.15$, seed~1) differs materially ($0.21\pct$ against $0.67\pct$) with
the same seed on a different GPU --- i.e.\ GPU-level non-determinism matters at the
${\sim}0.5$\,pp level on the lower branch. Wall-clock over the nine fresh points on this
GPU: baseline 7671\,s, 005 4148\,s, ratio 1.85$\times$ (branch solves only, CPU-pinned ---
explicitly \emph{not} comparable with the 2.72$\times$ full-mode figure).

\subsection{Item B --- building a correct Table~3 oracle}
\textbf{Pre-registered hypothesis.} On the semi-infinite domain the disturbance problem has
a continuous spectrum for $\alpha>K$; finite-domain solvers then produce roots that
accumulate as $L$ grows, and the lowest such root tends to the onset of oscillatory
far-field decay,
\begin{equation}
  \alpha^{*} = K + \frac{f_\infty^2}{4\,(1-\beta f_\infty^2)},
\end{equation}
rather than to a discrete eigenvalue. Convergence criteria, the support/refutation
conditions, and the rule that the PINN's stability stage is never re-tuned against the new
oracle were all fixed in advance.

\textbf{Deliverable.} \code{oracle-v2/table3\_oracle.py} (the bundled \code{oracle/} stays
byte-identical), sweeping three methods --- shooting with $g'(L)=0$, corrected Chebyshev,
and shooting with the exact asymptotic boundary condition --- over
$L\in\{10,12,16,20,24\}$, $S\in\{2.0,\ldots,2.4\}$ and both branches; ${\sim}25$ minutes on
CPU. Eq.~(31) was re-checked against the page-10 image: the coded operator is the exact
linearisation of steady Eq.~(13) with $+\alpha g'$, whereas the \emph{printed} equation drops
primes (it shows $+2ff''g_0$ and $-2ff''g_0$, which would cancel), so the linearisation is
taken as the authority.

\begin{table}[ht]
\centering
\caption{Table~3 oracle v2 ($\beta{=}0.1$, $K{=}0.1$, $A^{*}{=}0.1$, $\mathrm{Pr}{=}2$,
$B^{*}{=}0$).}
\label{tab:oraclev2}
\small
\begin{tabular}{@{}lrrrrr@{}}
\toprule
$S$ & \textbf{Lower $\alpha_1$ (converged)} & \textbf{Paper lower} & \textbf{Upper lowest finite-$L$ root} & $\alpha^{*}$ & \textbf{Paper upper} \\
\midrule
2.0 & $-0.49257$ & $-0.0376726$ & 0.7348 (converged) & 0.8205 & 0.8810518 \\
2.1 & $-0.61135$ & $-0.1290175$ & 1.0493 (converged) & 1.0882 & 1.1498579 \\
2.2 & $-0.70803$ & $-0.2006274$ & ${\sim}1.39$ ($0.36\pct$ drift $L20\!\to\!24$) & 1.4016 & 1.4609727 \\
2.3 & $-0.79412$ & $-0.2601093$ & not converged ($1.79\to1.92$) & 1.7828 & 1.8396721 \\
2.4 & $-0.87670$ & $-0.3115262$ & not converged ($2.28\to2.39$) & 2.2618 & 2.3116364 \\
\bottomrule
\end{tabular}
\end{table}

\paragraph{B1 --- the lower branch is converged, and Table~3's lower column does not
reproduce.} All three methods agree to $<10^{-7}$ relative at every $S$, and the value
changes by $<10^{-5}$ between $L{=}20$ and $L{=}24$; the pre-registered convergence
criterion is met at every row. The \emph{sign} agrees with the paper (lower branch
unstable), but the magnitudes are $2.8\times$ to $13\times$ larger than published. The
discrepancy is therefore on the paper's side, reproduced from the paper's own Eq.~(31)/(33).

\paragraph{B2 --- the upper-branch hypothesis was refuted, in a way the pre-registration did
not foresee.} At $S{=}2.0$ and $2.1$ the lowest finite-$L$ root does \emph{not} approach
$\alpha^{*}$ from above: it is converged ($<0.02\pct$ change over $L{=}16$--$24$) and lies
\emph{below} $\alpha^{*}$ ($-10\pct$, $-4\pct$). Higher roots do crowd toward $\alpha^{*}$ as
$L$ grows. The literal refutation condition (``converged and $>2\pct$ \emph{above}
$\alpha^{*}$'') was not met either --- the outcome was a third case, and this is stated as
such rather than being retrofitted into the pre-registered dichotomy. Interpretation offered
but not further tested: for $K<\alpha<\alpha^{*}$ both far-field modes decay monotonically,
and $g'(L)=0$ at large $L$ forces the slow-decaying mode's coefficient to vanish, so the
converged root is the ``fast-decay'' eigenvalue; under $g'(\infty)=0$ alone every
$\alpha>K$ still qualifies. At $S\geq2.2$ neither method converges (Chebyshev becomes
ill-conditioned for $L\gtrsim16$--$20$, producing spurious $|\alpha|\sim10^{2}$; shooting
drifts with $L$), so \textbf{no reliable upper magnitude is claimed for $S\geq2.2$}. What is
robust: every upper candidate root exceeds $K>0$, so the upper branch is stable --- the sign
agrees with the paper.

\paragraph{B3 --- attribution of the published upper column (informative only).} The bundled
\code{oracle/ref\_eig\_cheb.py} applies the wall boundary conditions at $\eta=L$ and the
far-field condition at the wall --- the rows are at swapped ends. At its own default
$L{=}12$ that placement yields upper values $0.88253$, $1.14325$, $1.45075$, $1.82675$,
$2.30169$: within $0.17$--$0.70\pct$ of \emph{all five} published Table~3 upper values. The
same placement gives positive values on the lower branch, so it does not reproduce the lower
column. This is recorded as \emph{consistent with, but not proof of}, the paper's upper
column having come from a computation with that boundary placement.

\paragraph{B4 --- energy-disturbance eigenvalues (new, Eq.~32 with $g_0=0$).} Never examined
before. All finite-$L$ energy eigenvalues are positive (lower branch $0.03$--$0.12$ at
$L{=}24$, decreasing with $L$), so they change neither sign conclusion. On the upper branch
at $S\geq2.3$ they fall below the momentum root at $L{=}24$ --- a further reason not to quote
an upper ``smallest eigenvalue'' magnitude.

\paragraph{Consequence for the PINN.} The PINN's lower-branch $\alpha_1$ at $S{=}2.0$
(baseline $-0.48820$, candidate-005 $-0.48808$) is within $0.9\pct$ of the converged
$-0.49257$. The evaluator's Table~3 magnitude gates compare against values this oracle does
not reproduce. The evaluator was left unedited (program-evolution contract); the oracle-v2
numbers are recorded as the recommended reference should a new evaluator version ever be
locked.

%=====================================================================
\section{Summary of findings}

\subsection{What reproduces}
\begin{itemize}[leftmargin=1.4em,itemsep=1pt]
  \item \textbf{Table 1 ($C_f$) and Table 2 ($\mathrm{Nu}$)} at the baseline row: upper
        $C_f$ to $0.004\pct$, lower $C_f$ to $0.42\pct$, $\mathrm{Nu}$ to $0.24$--$0.38\pct$
        --- the latter at the resolution of the independent oracle, which is itself
        $0.2$--$0.3\pct$ from the published values.
  \item \textbf{The critical suction parameter} $S_c = 1.849$ against the published
        $1.8391$ ($0.54\pct$).
  \item \textbf{The dual-solution structure and both stability signs}: upper branch stable,
        lower branch unstable, at every tabulated $S$ and by every method used.
  \item \textbf{Off-search-point behaviour}: three unseen Table~1 rows over three seeds,
        upper branch and both $\mathrm{Nu}$ within $0.34\pct$ of the converged oracle.
\end{itemize}

\subsection{What does not reproduce}
\begin{itemize}[leftmargin=1.4em,itemsep=1pt]
  \item \textbf{Table 3 eigenvalue magnitudes.} The lower column is off by factors of
        $2.8$--$13$; three independent solvers of the paper's own Eq.~(31) agree with each
        other to $<10^{-7}$ relative and not with the paper. The upper column is
        reproducible to $0.17$--$0.70\pct$ only by a boundary-row placement that is
        mathematically wrong, and the upper magnitude is in any case not identifiable
        either by this PINN (continuous spectrum for $\alpha>K$) or reliably by the oracle
        at $S\geq2.2$.
  \item \textbf{The signs in Table~2.} The article prints negative values under the heading
        $\mathrm{Nu}\,\mathrm{Re}_x^{-1/2}$ while defining that quantity as $-\theta'(0)>0$;
        both the PINN and the independent oracle return positive values of the same
        magnitude. Verified against page~7; see Section~\ref{sec:beat}.
  \item \textbf{Candidate-003's claimed accuracy gain.} A real CPU-level improvement
        ($0.19\pct\to0.018\pct$ lower-$C_f$ error at a 300-iteration cap) that did not
        survive full-mode evaluation.
  \item \textbf{The session-1 CPU full-mode run.} Never completed; the output file is empty
        and is documented as empty.
\end{itemize}

\subsection{What was improved}
\begin{itemize}[leftmargin=1.4em,itemsep=1pt]
  \item \textbf{2.72$\times$ end-to-end same-GPU speed-up} (1389.9\,s $\to$ 511.6\,s) from
        candidate-005, with accuracy unchanged at the search point and
        \emph{equivalent} --- under a pre-registered rule --- on unseen rows and seeds.
        Stage-level: $2.48\times$ (Table~1/2), $2.45\times$ (fold), $3.13\times$
        (stability).
  \item The dominant mechanism is \textbf{replacing three nested autograd passes with
        closed-form Taylor-mode derivative jets}, verified gradient-identical to
        $2\times10^{-14}$ before evaluation; the rest comes from adaptive LBFGS stopping and
        a cheaper $L_{\mathrm{MAP}}$ pilot.
  \item A \textbf{corrected, converged Table~3 oracle} now exists where none did, including
        the previously unexamined energy-disturbance spectrum.
\end{itemize}

%=====================================================================
\section{Did this work beat the article?}
\label{sec:beat}

The honest verdict is \textbf{partial, and the partition is sharp}: this work
\emph{supersedes} the article on its stability table and provides a better numerical
reference for its flow tables, while the article remains ahead on raw computational cost
and on breadth of parameter coverage. The two claims should not be traded against each
other, and the speed figure reported in Section~6 is \emph{not} a claim against the paper's
method.

\begin{table}[ht]
\centering
\caption{Axis-by-axis comparison against Pandey \textit{et al.}}
\label{tab:beat}
\small
\begin{tabular}{@{}p{3.9cm}p{3.7cm}p{4.3cm}p{3.0cm}@{}}
\toprule
\textbf{Axis} & \textbf{The article} & \textbf{This work} & \textbf{Verdict} \\
\midrule
Stability eigenvalue magnitudes (Table~3)
  & Published to 7 decimals; not recoverable from its own Eq.~(31)
  & Converged values by three independent methods, agreeing to $<10^{-7}$
  & \textbf{superseded} \\
Energy-disturbance spectrum (Eq.~32)
  & Not examined
  & Computed; all eigenvalues positive, no sign change
  & \textbf{new} \\
Far-field reference for Tables~1/2
  & $\eta_\infty = 10$ truncation
  & Converged $\eta_\infty = 40$ oracle
  & \textbf{improved} \\
Table~1/2 values themselves
  & 8 significant figures
  & PINN within $0.004$--$0.42\pct$
  & matched, not beaten \\
Internal consistency of Table~2
  & Signs contradict the paper's own definition
  & Inconsistency identified and verified
  & \textbf{defect found} \\
Compute per parameter point
  & RK4 $+$ secant: milliseconds
  & 511.6\,s for the full three-stage pipeline
  & \textbf{article wins} \\
Parameter coverage and physical interpretation
  & Dense grids, 14 figures, Section-6 synthesis
  & One gated row, $S_c$, $S{=}2.0$ stability, 3 unseen rows
  & \textbf{article wins} \\
Auditability of the record
  & Printed tables
  & Hashes, locks, pre-registrations, raw sweeps, negative results
  & \textbf{this work} \\
\bottomrule
\end{tabular}
\end{table}

\subsection{Where this work supersedes the article}
Table~3 is the one place where the article's published numbers are shown to be wrong rather
than merely imprecise. Three independent solves of the paper's \emph{own} Eq.~(31) --- IVP
shooting with $g'(L)=0$, a corrected Chebyshev discretisation, and shooting with the exact
asymptotic boundary condition --- agree with one another to $<10^{-7}$ relative at every
tabulated $S$ and change by $<10^{-5}$ between $L{=}20$ and $L{=}24$, while differing from
the published lower-branch column by factors of $2.8$ to $13$ (Table~\ref{tab:oraclev2}).
No correct solver can land within $1\pct$ of $-0.0376726$ at $S{=}2.0$. The qualitative
conclusion the article draws from Table~3 --- upper branch stable, lower branch unstable at
every $S$ --- is unaffected: every method reproduces both signs. What does not survive is
the quantitative column.

Two further contributions have no counterpart in the article: a mechanical candidate
explanation for how the published \emph{upper} column could have arisen (a boundary-row
placement that is mathematically wrong but reproduces all five values to within $0.70\pct$),
and the first computation of the energy-disturbance spectrum of Eq.~(32), which turns out to
be everywhere positive and therefore to change neither sign conclusion --- while also
showing that on the upper branch at $S\geq2.3$ it falls \emph{below} the momentum root, one
more reason the article's upper magnitudes should not be quoted.

\subsection{Where this work improves the article's numbers without overturning them}
The article truncates the far field at $\eta_\infty = 10$. The converged $\eta_\infty = 40$
oracle built here shows what that costs on the slowly-decaying lower branch: the published
lower $C_f$ is $0.17\pct$ below the converged value at $S{=}1.95$ and $0.78\pct$ below it at
$S{=}2.1$, with the gap growing as the branch decays more slowly. On $\mathrm{Nu}$ the
oracle sits $0.20\pct$ (upper) and $0.29\pct$ (lower) from the published values. The
practical consequence is that the article's own tables --- not the PINN --- are the limiting
reference: the PINN's $0.24$--$0.38\pct$ agreement on $\mathrm{Nu}$ is at the resolution of
the best available independent solve, so no finer comparison is meaningful. The PINN itself
is therefore a faithful reproduction, not an improvement, of Tables~1 and~2; the
\emph{oracle} is the improvement.

\subsection{A verified internal inconsistency in Table~2}
The article defines $\mathrm{Nu}\,\mathrm{Re}_x^{-1/2} = -\theta'(0)$ (Section~3), and with
$\theta(0)=1$, $\theta(\infty)=0$ this is positive for the cooling wall considered. Page~7,
Table~2, nevertheless prints \emph{negative} values under the heading
$\mathrm{Nu}\,\mathrm{Re}_x^{-1/2}$: the baseline row
($S{=}1.95$, $K{=}0.1$, $\beta{=}0.1$, $\mathrm{Pr}{=}2$, $A^{*}{=}B^{*}{=}0.1$) reads
$-3.35807824$ and $-3.14427924$, and every row in the table carries the same sign.
Both the PINN and the independent \code{scipy.solve\_bvp} oracle return $+3.350$ and
$+3.132$ --- magnitudes that agree to $0.2$--$0.4\pct$. The table therefore appears to
tabulate $\theta'(0)$ while labelling it $-\theta'(0)$; it is a presentational slip, not a
numerical error, and the figures on the same page (Figs.~9--10, plotted on positive axes)
are consistent with the positive convention. It is recorded here because a reader copying
Table~2 verbatim would carry the wrong sign into downstream work. The evaluator gates on
magnitude, so nothing in this project is affected. This was verified directly against
page~7 of the archived copy whose SHA-256 is recorded in Table~\ref{tab:prov}.

\subsection{Where the article still wins}
\begin{itemize}[leftmargin=1.4em,itemsep=2pt]
  \item \textbf{Computational cost.} The article's shooting method (RK4, $\Delta\eta=0.01$,
        secant to $10^{-5}$) solves a branch in milliseconds. The fastest program produced
        here needs 511.6\,s for the full three-stage pipeline. The $2.72\times$ reported in
        Section~6 is measured against the supplied PINN baseline on identical hardware and
        says nothing about the paper's own solver; a PINN was never going to win that
        comparison, and nothing in this record claims it does. The case for the PINN is
        that it imposes the far-field conditions exactly, passes smoothly through the
        turning point, and provides a third independent route to the same claims --- not
        that it is cheap.
  \item \textbf{Breadth.} The article sweeps dense grids in $S$, $K$, $\beta$,
        $\mathrm{Pr}$, $A^{*}$ and $B^{*}$, with fourteen figures and a synthesis of the
        physical mechanisms. This project gates one parameter row, plus $S_c$, plus the
        $S{=}2.0$ stability row, with three unseen rows over three seeds as confirmation.
        The qualitative conclusions of the article's Section~6 were transcribed as
        constraints to protect, never independently re-derived across the full grid.
  \item \textbf{Physical interpretation.} No claim about the physics of the Maxwell
        shrinking-sheet problem originates here.
\end{itemize}

\subsection{What would actually constitute beating it}
Two things, neither yet done:

\begin{enumerate}[leftmargin=1.4em,itemsep=2pt]
  \item \textbf{Publish the Table~3 correction.} The evidence is complete and preserved ---
        three agreeing methods, domain-length studies, pre-registered criteria applied
        mechanically, and a candidate attribution for the published upper column. What is
        missing is the write-up and the courtesy of putting it to the authors; the article
        is an uncorrected proof, so the window for a correction may still be open.
  \item \textbf{Execute the D-track} (Section~4.3, \code{candidate-mechanisms.md}). A single
        parametric network $f(\eta; S,\beta,K)$, $\theta(\eta; S,\beta,K,\mathrm{Pr},A^{*},B^{*})$
        would reproduce the entire Table~1/2 sweep from one training run and extend past it
        --- including the non-existence region below $S_c$, which the article maps only by
        where its solver stops converging. Differentiating that network with respect to $S$
        gives continuous fold-tracking and a continuous $\alpha_1(S)$, which would verify as
        \emph{one} claim that the turning point and the stability-exchange point coincide.
        The article asserts this through two separate tables and never checks it. That is
        the result that would be beyond the paper rather than alongside it, and it needs
        held-out $(S,\beta,K)$ points and its own evaluator, not a candidate folder.
\end{enumerate}

%=====================================================================
\section{Methodological discipline actually followed}

This is a large part of the value of the record and is worth stating explicitly:

\begin{enumerate}[leftmargin=1.4em,itemsep=2pt]
  \item \textbf{Pre-registration before results.} Three pre-registration documents
        (\code{GPU-SEARCH-PREREGISTRATION.md}, \code{OPEN-ITEMS-PREREGISTRATION.md}, and the
        \code{oracle-v2} criteria mechanically applied by \code{judge.py} into
        \code{judge.md}) fix acceptance rules, thresholds and convergence criteria in
        advance. When an outcome fell outside the foreseen cases (Item~B upper branch), that
        is reported as such rather than reinterpreted.
  \item \textbf{The evaluator was never edited}, not even when two of its gates were proven
        unmeetable by any correct solver. The gates are reported, excluded from acceptance
        with a documented reason, and the corrected reference values are recorded for a
        future locked evaluator version.
  \item \textbf{No cross-hardware timing claims.} CPU and GPU numbers are reported
        separately and never summed; session-3 timings are never compared with session~2's
        because the GPU differs; the $1.85\times$ fresh-row ratio is explicitly marked as
        not comparable with the $2.72\times$ full-mode figure.
  \item \textbf{Hash-verified reconstruction.} The winner was rebuilt from the locked
        baseline plus its recorded patch and verified byte-identical before confirmation.
  \item \textbf{Paired, idle-machine, serial measurement} with recorded timelines,
        \code{nvidia-smi} output, environment dumps and SHA-256 manifests next to every run.
  \item \textbf{Smoke runs are never reported as reproduction.} Stated in the evaluator
        docstring, the project notes, and the skill.
  \item \textbf{Negative and null results are kept} --- candidate-003's non-reproducing
        gain, candidate-004's within-noise speed-up, the dead CPU run, the refuted
        spectrum hypothesis.
  \item \textbf{Honest statistical limits.} ``Equivalent'' is defined as ``no gross
        difference detected at $n=9$'', not as equality; the non-pairing of seeds across the
        two programs is noted.
  \item \textbf{Source material untouched.} Zero modified files against upstream; all new
        work in two clearly separated directories.
\end{enumerate}

%=====================================================================
\section{Evidence inventory}

\begin{longtable}{@{}>{\raggedright\arraybackslash}p{7.0cm}>{\raggedright\arraybackslash}p{8.7cm}@{}}
\caption{Principal artifacts under
\code{initiatives/maxwell-shrinking-pinn/}
\code{projects/pinn-reproduction/}.}\label{tab:inventory}\\
\toprule
\textbf{Path} & \textbf{Content} \\
\midrule
\endfirsthead
\toprule
\textbf{Path} & \textbf{Content} \\
\midrule
\endhead
\code{../BRIEF.md} & Original request, provenance hashes, desired outcome, constraints \\
\code{NOTES.md} & Methodology, per-session status, limitations, reproduction commands \\
\code{paper/paper-summary.md} & Transcribed equations, BCs, Tables 1--3 with gate markings, Section-6 conclusions \\
\code{paper/}\allowbreak\code{pagination\_FINMEC\_100352-1.pdf} & Read-only copy of the source PDF \\
\code{baseline/pinn\_maxwell\_shrinking.py} & Unmodified three-stage PINN (the program to evolve) \\
\code{oracle/reference\_bvp.py} & Independent \code{scipy.solve\_bvp} ground truth (unmodified) \\
\code{oracle/ref\_eig\_cheb.py} & Independent Chebyshev eigenvalue solver (unmodified; BC rows at swapped ends) \\
\code{evaluator.py} & Candidate-proof evaluator; smoke/full modes, \code{--fold}, \code{--stability}, \code{--oracle} \\
\code{baseline-lock/}, \code{baseline-lock-gpu/} & Hash-locked baseline $+$ evaluator snapshots, CPU and GPU \\
\code{candidates/candidate-001..005/} & Per-hypothesis folder: \code{NOTE.md}, \code{patch.diff}, full source, \code{smoke.json}, \code{full.json}, stderr logs \\
\code{candidates/LINEAGE.md} & Lineage table with timings, speed-ups and gate outcomes \\
\code{runs/full\_confirmation\_gpu.json} & Same-GPU full-mode reference (1451.8\,s) \\
\code{runs/GPU-SEARCH-PREREGISTRATION.md} & Acceptance rule and the two non-identifiability arguments \\
\code{runs/GPU-SEARCH-RESULTS.md} & Winner confirmation, paired timings, fresh-parameter check, verdict \\
\code{runs/confirm-005/} & Paired baseline/winner runs, fresh-seed run, \code{env.txt}, \code{nvidia-smi.txt}, \code{timeline.txt}, run script \\
\code{runs/OPEN-ITEMS-PREREGISTRATION.md} & Session-3 rules for Items A and B, with program hashes \\
\code{runs/fresh-paired-2060s/ANALYSIS.md} & Item A result: per-point table, means/maxima, verdict, cross-GPU note \\
\code{runs/fresh-paired-2060s/analyze.py} & The analysis script that produced it \\
\code{runs/full\_table12\_confirmation.json} & Session-1 CPU run that never completed (empty, documented) \\
\code{confirm\_fresh.py} & Unseen-row, multi-seed confirmation harness \\
\code{diagnostics/eig\_shooting.*}, \code{eig\_cheb\_fixed.*} & Independent eigenvalue solves with domain-length studies (code $+$ output) \\
\code{oracle-v2/} & Corrected Table~3 oracle: solver, raw JSON, \code{judge.py}/\code{judge.md}, \code{README.md} with the full table \\
\code{../../../skill-lab/}\allowbreak\code{evolve-maxwell-pinn-gpu/} & \code{SKILL.md} and \code{references/candidate-mechanisms.md} (14 mechanisms, A--D) \\
\bottomrule
\end{longtable}

\subsection{Commit history}
\begin{table}[ht]
\centering
\small
\begin{tabular}{@{}llp{9.6cm}@{}}
\toprule
\textbf{Commit} & \textbf{Date} & \textbf{Subject} \\
\midrule
\code{9ee93bc} & 2026-11-05 & Maxwell shrinking-sheet PINN: Open Discovery GPU reproduction and candidate search \\
\code{d5fd1ba} & 2026-11-05 & Confirm candidate-005: 2.72$\times$ same-GPU speed-up; record results and lineage \\
\code{2e1be62} & 2026-11-06 & Session 3: pre-register open items; Table~3 oracle v2; start paired fresh-seed runs \\
\code{68b53ec} & 2026-11-06 & Sync: paired fresh-seed run progress (baseline seeds on 2060 SUPER) \\
\code{4385559} & 2026-11-06 & Sync: candidate-005 fresh-seed run progress (2060 SUPER) \\
\code{30c62ea} & 2026-11-06 & Close open items: baseline vs 005 fresh-seed spread equivalent; Table~3 oracle v2 results \\
\bottomrule
\end{tabular}
\end{table}

\noindent The 6.7\,GB GPU venv is excluded from version control; \code{open-discovery}'s own
upstream history is preserved as \code{open-discovery/.git-upstream} (rename back to
\code{.git} to restore it as a standalone checkout).

%=====================================================================
\section{How to reproduce}

\begin{verbatim}
# environment (GPU; keeps sm_61 kernels for a GTX 1080 Ti)
uv pip install -r amit-pinn-article/requirements-gpu-cu126.txt \
  --extra-index-url https://download.pytorch.org/whl/cu126

cd amit-pinn-article/open-discovery/initiatives/maxwell-shrinking-pinn/\
projects/pinn-reproduction

# cheap pipeline sanity gate (~15 s GPU, ~127 s CPU)
python evaluator.py baseline/pinn_maxwell_shrinking.py

# the actual paper-accuracy claim (~1450 s on a 1080 Ti for all three stages)
python evaluator.py baseline/pinn_maxwell_shrinking.py \
  --mode full --oracle --fold --stability

# unseen-row, multi-seed confirmation
python confirm_fresh.py baseline/pinn_maxwell_shrinking.py 1 2 3

# rebuild the winner from the locked baseline + its recorded patch
cp baseline-lock-gpu/baseline/pinn_maxwell_shrinking.py /tmp/c005.py
patch /tmp/c005.py < candidates/candidate-005/patch.diff

# corrected Table 3 oracle (~25 min, CPU)
python oracle-v2/table3_oracle.py && python oracle-v2/judge.py

# independent eigenvalue diagnostics
python diagnostics/eig_shooting.py
python diagnostics/eig_cheb_fixed.py

# re-lock after any intentional change to baseline/ or evaluator.py
python ../../../../program-evolution/scripts/lock_baseline.py \
  --baseline baseline/pinn_maxwell_shrinking.py --evaluator evaluator.py \
  --goal "<state the goal>" --output baseline-lock-v2 --timeout-seconds 1800
\end{verbatim}

%=====================================================================
\section{Status, limitations and possible next steps}

\paragraph{Status.} \textbf{Reproduction complete to the extent the paper allows}, with no
open items. Tables~1 and~2 and $S_c$ reproduce; Table~3 signs reproduce and magnitudes do
not, with the discrepancy localised to the paper side and evidenced by three independent
solvers. Candidate-005 is confirmed $2.72\times$ faster on the same GPU and
accuracy-equivalent to the baseline.

\paragraph{Limitations to carry forward.}
\begin{itemize}[leftmargin=1.4em,itemsep=1pt]
  \item The evaluator still carries two Table~3 magnitude gates that no correct solver can
        pass, so \code{valid=false} on \emph{every} full-mode run including the untouched
        baseline. Any future locked evaluator version should adopt the oracle-v2 reference
        values.
  \item The lower-branch $C_f$ shows a ${\sim}0.03$--$1.0\pct$ seed spread against the
        converged oracle. This is inherent to the baseline architecture, not to any
        candidate, and is partly GPU-level non-determinism (${\sim}0.5$\,pp at identical
        seeds across devices).
  \item No reliable upper-branch $\alpha_1$ magnitude exists at $S\geq2.2$ by any method
        used here.
  \item All accuracy gates are at a single parameter point; $n=9$ off-search-point checks
        detect only gross differences.
  \item Mechanisms A1 (batching both branches via \code{vmap}), A4 (Adam/LBFGS rebalance),
        A5 (vectorising the fold continuation loop) and A6 (larger collocation batches)
        from the mechanism menu remain untried.
\end{itemize}

\paragraph{Possible next steps.} The D-track in
\code{references/candidate-mechanisms.md} is genuinely new science beyond the paper and
would require a \emph{new project} with its own evaluator and held-out $(S,\beta,K)$ grid
points: a parametric/meta-PINN $f(\eta;S,\beta,K)$ reproducing the entire Table~1/2 sweep in
one training run (and mapping the non-existence region below $S_c$); continuous
fold-tracking by differentiating that network with respect to $S$; parametric stability
producing a continuous $\alpha_1(S)$ and locating $\alpha_1=0$ directly, which would check
as \emph{one continuous claim} that the turning point and the stability-exchange point
coincide --- something the paper asserts through separate tables but never verifies; and a
small multi-seed ensemble reporting $C_f$/$\mathrm{Nu}$ spread as a robustness companion.

\vspace{1em}
\noindent\rule{\textwidth}{0.4pt}

\paragraph{A note on scope.} This initiative is a reproduction plus
harness-and-optimisation study. It asserts no new claim about the physics of the Maxwell
shrinking-sheet problem. The one finding with implications beyond the harness --- that
Table~3's eigenvalue magnitudes are not recoverable from the paper's own Eq.~(31), while a
specific incorrect boundary-row placement does recover its upper column to within
$0.7\pct$ --- is stated as evidence, with the raw sweeps, three independent methods, and
pre-registered criteria preserved so a reader can judge it independently.

\end{document}
